\documentclass[a4,25pt]{article}

\usepackage{indentfirst}
\usepackage{latexsym}
\usepackage[empty]{fullpage}
\usepackage{titlesec}
\usepackage{marvosym}
\usepackage[usenames,dvipsnames]{color}
\usepackage{verbatim}
\usepackage[hidelinks]{hyperref}
\usepackage[utf8]{inputenc}
\usepackage{fancyhdr}
\usepackage{multicol}
\usepackage{hyperref}
\usepackage{csquotes}
\usepackage{tabularx}
\hypersetup{colorlinks=true,urlcolor=black}
\usepackage[11pt]{moresize}
\usepackage{setspace}
\usepackage{fontspec}
\usepackage[inline]{enumitem}
\usepackage{array}
\newcolumntype{P}[1]{>{\centering\arraybackslash}p{#1}}
\usepackage{anyfontsize}
\usepackage{fontspec}
\usepackage{academicons}
\usepackage{fancyhdr}
\usepackage{fontawesome5}
\usepackage{indentfirst}
\usepackage{lastpage}
\usepackage{parskip}[0.5cm]
\usepackage{hyperref}
\usepackage{url}
%%%% Set Margins
\usepackage[margin=1.3cm, top=1.3cm, bottom=1.5cm]{geometry}
\usepackage[english]{babel}
\usepackage[document]{ragged2e}

%%%% Set Fonts
\setmainfont[
BoldFont=EBGaramond-Bold.ttf, %SourceSansPro-Bold.otf
ItalicFont=EBGaramond-Italic.ttf
]{EBGaramond-Regular.ttf}
\setsansfont{EBGaramond-SemiBold.ttf}

%%%% Set Page
\pagestyle{fancy}
\fancyhf{} 
\fancyfoot{}
\renewcommand{\headrulewidth}{0pt}
\renewcommand{\footrulewidth}{0pt}

%%%% Set URL Style
\urlstyle{same}

%%%% Set Indentation
\raggedbottom
\raggedright
\setlength{\tabcolsep}{0in}

%%%% Set Secondary Color
\definecolor{UI_blue}{RGB}{0,0,0}

%%%% Define New Commands


%%%% Set Sections formatting
\titleformat{\section}{
 \color{UI_blue} \scshape \raggedright \large \bfseries
}{}{17.3em}{}[\vspace{-10pt} \hrulefill \vspace{-2pt}]

%%%% Set Subtext Formatting
\newcommand{\subtext}[1]{
#1\par\vspace{-0.2cm}}

\newcommand{\subtextit}[1]{\vspace{0.15cm}
 \textit{ #1 \vspace{-0.2cm}} }
 %%%
\newcommand{\cps}{{\color{red} \quad \textbf{|}\quad}}
\newcommand{\cpshalf}{{\color{red} \hspace{0.2em}\textbf{|}\hspace{0.2em}}}
%%%% Set Item Spacing
\setlist[itemize]{align=parleft,left=10pt..2em}

%%%% New Itemize "Zitemize" Formatting - tighter spacing than itemize
\newenvironment{zitemize}{
\begin{itemize}\itemsep0pt \parskip0pt \parsep1pt}
{\end{itemize}\vspace{-0.5cm}}


%%%% Define Skills Bold Formatting
\newcommand{\hskills}[1]{
\textbf{\bfseries #1} }

%%%% Set Subsection formatting
\titleformat{\subsection}{\vspace{-0.1cm} 
\bfseries \fontsize{11pt}{2cm}}{}{0em}{}[\vspace{-0.2cm}]
%%%%%%

%%%%%%% --------------------------------------------------------------------------------------
\rfoot{{\par\fontfamily{ttf} \selectfont \footnotesize \emph{ \thepage \hspace{1pt} of \pageref{LastPage}}}}
\setlength{\parindent}{4em}
\setlength{\parskip}{0.5cm}
%%%%%%% --------------------------------------------------------------------------------------

\def\Vrulefill#1#2{
    \vspace{-10pt}
    \leavevmode%
    %\hskip-.2in%
    \leaders%
    \vtop{\hsize=.0025in\vskip#1#2}%
    \hfill%
    %\hskip.3in%
    \vspace{-2pt}
}


\begin{document}

%%%%%%% --------------------------------------------------------------------------------------
%%%%%%%  HEADER
%%%%%%% --------------------------------------------------------------------------------------
{{\huge \scshape Statement of Purpose}}\\
\large{Arman Asgharpoor Golroudbari} \hfill { \small \textbf{10 December 2022}} \\

\href{https://ArmanAsq.github.io}{\textbf{\small{https://ArmanAsq.github.io}}}
\hfill {Aeronautics and Astronautics}\\

\Vrulefill{-.05in}{-}
%\rule{10cm}{0.1em}
%%%%%%% --------------------------------------------------------------------------------------
%%%%%%%  Research Interest
%%%%%%% --------------------------------------------------------------------------------------
\justifying

\begin{large} 
Dear Graduate Admissions, \par
\vspace{-0.25cm} \quad 
I am an avid astronomer passionate about exploring the night sky. This is why I investigate how AI can solve technical issues surrounding Space Exploration. I want an interdisciplinary Ph.D. in emerging fields at the intersection of astronomy, machine learning, and space flight because I want to overcome technical problems associated with navigation in interstellar missions. As a result, I am applying to Harvard University's Department of Astronomy for a Ph.D. I hope to get the opportunity to work with Professor Douglas Finkbeiner, among your other faculty members, in these areas. \par
%\subsection{Education}
Reading \textit{Einstein}'s book \textit{The Meaning Of Relativity} sparked my interest in astronomy and cosmology when I was sixteen. This passion led me to Amateur Astronomy, where I study the night sky with my 10" Newtonian telescope. My fascination with the universe and the mysteries of the cosmos motivates me to participate in various astronomy events and to become a member of the Night Sky Institute. Consequently, I began to study celestial objects, first the Orion Nebula and then the Andromeda Galaxy. My eyes were blown away by the sight of a giant cloud of dust and gas in space. But none of these experiences satisfied me. There has always been a fascination for me with the origin and evolution of the universe. I'm not only interested in observing the cosmos from behind a telescope but also in exploring it. To make my dream come true, I chose to study Space Engineering as my major in graduate school.

\vspace{-0.25cm}\quad As a result of my background in Avionics Engineering and passion for space, I decided to pursue a master's degree in Space Engineering, focusing on integrating AI into navigation. As part of the Spacecraft System Design course, I used European Cooperation for Space Standardization (ECSS) to conceptualize a FireSat mission based on the \textit{Space Mission Analysis and Design} book. This project sparked my interest in space missions and Interstellar Objects, particularly the concept of sending a spacecraft to the local interstellar medium, such as Breakthrough Starshot. Subsequently, I took various relevant courses and taught myself deep learning because I saw its strong potential as a vision-based navigation method. I also participated in Oxford Machine Learning Summer School to gain in-depth knowledge of various AI techniques. It allows me to examine vision-based navigation techniques and implement path-planning algorithms using Python in the ROS environment. \par
\vspace{-0.25cm}\quad Because I want to become an entrepreneur in the space sector, I enrolled in an interdisciplinary MBA. Through various coursework and extracurricular activities, I learned a great deal about what it takes to start and run a successful business, from developing a business plan and marketing strategy to financial management and fundraising. My studies and internships enabled me to network with people in the industry to learn from their experiences. Studying Space Engineering alongside an interdisciplinary MBA helped me develop leadership and teamwork abilities. To enhance my leadership skills and communicate effectively with students, I spent more than three years working as a thesis supervisor and university lecturer. I have supervised undergraduates on their thesis projects and taught eleven undergraduate courses at the Aviation Industry Training Center. Additionally, I have served as a martial arts instructor since 2016, teaching students from various backgrounds. These experiences encouraged me to continue my career in academia. \par
%\subsection{Research Experience}
\vspace{-0.25cm}\quad 
While pursuing my MSc in Space Engineering, I conducted navigation-based research projects designed to develop deep learning models to solve the attitude estimation problem. My master's thesis focuses on the application of AI in navigation. I initially developed inertial navigation frameworks with TensorFlow and PyTorch using the Random Search algorithm. Then I tuned them with optimization techniques, including Population-Based Training and Bayesian Optimization. As part of my master's thesis, I developed the first end-to-end inertial attitude and heading estimation model based on the Magnetometer+IMU measurements. These projects enabled me to learn the advanced level of deep learning algorithms as well as programming and problem-solving skills.\par
%\subsection{Research Interest}
\vspace{-0.25cm}\quad 
I want to pursue graduate studies in computer vision and AI to address the issue of interstellar missions navigation (e.g., solar sailing). I am fascinated by the idea of using deep learning for interstellar missions. 
Due to relativistic speeds, solar sailing missions such as Breakthrough Starshot will experience time dilation, length contraction, and relativistic Doppler shift as they travel at a fraction of the speed of light. These effects lead to calculation errors. Therefore, popular methods such as AutoNav \cite{team2000autonomous} and XNAV \cite{graven2008xnav}, as well as Newtonian mechanics algorithms, would be inaccurate in these missions. Therefore, an autonomous navigation model must be developed that incorporates Newtonian and Relativistic mechanics in missions to interstellar medium \cite{yucalan2021autonomous,yucalan2021optimal,yucalan2022autonomous}. Also, in the Solar System, stars may appear to be pointing in a fixed direction in space with only minor perturbations caused by stellar aberrations, parallax, and relative motion, whereas spacecraft traveling outside the solar system observes stars to be located differently. In order to solve these problems, I envision developing a 3D star catalog, or to put it in a wider perspective, a 3D map of the observed universe (mainly the Milky Way galaxy and Pulsars) using the Gaia dataset. As there are currently no comprehensive 3D star catalogs, and most extensive catalogs only cover a fraction of Milky Way Galaxy stars, many of their distances are uncertain by light-years\cite{mckee2022navigation}. Using this map, I want to train deep learning models that could solve the localization problem in interstellar missions using Physics-Informed Neural Networks that incorporate Newtonian and relativistic mechanics. However, I am also looking forward to exploring new ideas. \par
%\subsection{Conclusion}
\vspace{-0.25cm}\quad As a prospective Ph.D. student with a strong background in AI and a passion for space-related research, I believe that the Ph.D. program in Harvard University's Department of Astronomy is an ideal opportunity for me to have concrete research experience and craft viable career plans. Further, I want to pursue a Ph.D. because it is a perfect chance to work on cutting-edge projects in collaboration with renowned scientists. It is a privilege to study in this department under the guidance of its distinguished faculty members. Also, I am fascinated by the many exciting research projects taking place in this department. In particular, I am interested in the research that was conducted by \textbf{Professor Douglas Finkbeiner} to build a 3D map of dust \cite{green20193d} using Gaia, Pan-STARRS 1, and 2MASS. In my studies, I want to learn about all aspects of Astronomy and develop my research skills. With my motivation, intellectual ability, and preparation, I am confident I can embark on this exciting and challenging path and make significant original contributions to you and your research project.\par
\vspace{-0.25cm} \quad As a highly motivated student, I demonstrated a solid academic record and hands-on experience in various research areas, which also manifested in my GPA of 4.0/4.0 in both my MBA and MSc programs. Given that my research interests and area of expertise align well with machine learning and astronomy, I hope to get the opportunity to work with Professor Douglas Finkbeiner, Daniel Eisenstein, or Benjamin Johnson, among your other faculty members in these areas. \par
\vspace{-0.25cm} \quad Harvard University is one of the world's leading centers for research and innovation, making it an ideal place to pursue my Ph.D. \par
I look forward to being given the opportunity to join Harvard University as a graduate student in your esteemed department.\par


In advance, I appreciate your time and consideration. \par
\vspace{-0.4cm} Sincerely Yours, \par
\vspace{-0.4cm}Arman Asgharpoor Golroudbari \par

\end{large}
\bibliographystyle{unsrt}
\small
\bibliography{refs}

\end{document}